\paragraph{Bounding the incidence form by an integral section.}
We detail the Chow estimate used above. The integral lower bound
needed is at one specified multiple of the light weight. Let
\(M=(s+2)N\), \(t=1+\lceil P_O\rceil\), and suppose \(\delta>0\).
Put
\[
 h=8M^2(s+1)\delta t,\qquad h'=8hM(s+1)\delta .
\]
For every component \(Z\) the input is
\begin{equation}\label{eq:chow-lower-section}
 \operatorname{mean}_{\sigma}\log\|g|_Z\|_{\sup,\sigma}
                    \ge-h P_0 S_L
\end{equation}
whenever \(g\) is an integral linear combination of products
of \(h\) light columns and \(g|_Z\ne0\). A lower bound on the
supremum over the union of all components alone would not
give this inequality. The integer \(h\) is fixed before applying it.

\begin{proposition}[The component Chow form]
\label{prop:integral-chow-bound}
With the componentwise input \eqref{eq:chow-lower-section} at the
chosen \(h\), the incidence forms of the selected family satisfy
\[
 \sum_Z h_0(F_Z)
 \le(s+2)\delta(P_0 S_L+G_c S_L+\log N+1).
\]
The sum is over the same components as the slice degree.
\end{proposition}

\begin{proof}
\proofstep{The incidence polynomial under a linear substitution.}
Write \(Y_i=\sum_l U_{il}w_l\), \(0\le i\le s+1\), with
all \(M\) entries of \(U\) independent. Let \(q(U,Y)\) be
the primitive irreducible equation of the projected image of
the given light component in \(\mathbb P^{s+1}_{K(U)}\).
Set \(d_0=\deg_Y q\), \(u=\deg_U q\), and let \(m_Z\)
be the covering multiplicity used in \(F_Z\). We use the
incidence construction with
\[
 1\le m_Zd_0\le\delta_Z,\qquad u\le(s+1)\delta_Z,\qquad
                    h_0(F_Z)\le m_Z h_{\mathbb T}(q).
\]
Here \(h_{\mathbb T}\) uses the same Gauss norms and the
mean log absolute value on unit \emph{circles} in all
variables. To see the incidence comparison, let \(F_0\) be
the Chow polynomial of the light image, so \(F_Z\) is its
\(m_Z\)-th power up to scalar.

On \(Y_0\ne0\) substitute in the \(s+1\) rows of the Chow polynomial the incidence coordinates
\[
  U'_{i l}=Y_0U_{i+1,l}-Y_{i+1}U_{0l}.
\]
After inverting \(Y_0\) this is a change of independent row
variables followed by multiplication by units. The resulting
equation is irreducible up to powers of \(Y_0\); removing
those powers gives \(q\). Geometrically it tests equality
of a projected image point and \([Y_0:\cdots:Y_{s+1}]\).
For a general line in the target its inverse image is cut by
\(s\) light hyperplanes. Counting those points bounds
\(m_Zd_0\) by \(\delta_Z\); the row substitution gives the
stated bound for \(u\). The homogeneous equations vanishing
after 
\[
  Y_i=\sum U_{il} w_l
\]
 form the principal prime \((q)\):
over \(K(U)\) the image has dimension \(s\) (a general linear
projection preserves dimension), and primitivity descends
divisibility to \(K[U,Y]\). Equivalently the differential
argument for the generic rows proves this dimension statement.

In the substitution into \(F_Z\) the coefficient of the
highest power of \(Y_0\) is the original polynomial in the
other rows of \(U\). Its log Gauss norm and torus mean do not
exceed those of the substituted polynomial, by extraction
of a leading coefficient and Jensen on the \(Y_0\) circle.
In particular for \(P(Y_0)=\sum_{j=0}^d a_jY_0^j\) the local coefficient estimate is
\[
  \|a_d\|_{\rm Gauss}\le\max_{0\le j\le d}\|a_j\|_{\rm Gauss}=\|P\|_{\rm Gauss}.
\]
The mean on spheres of \(F_Z\) is at most its torus mean:
condition on the radii and apply monotonicity of circle
means of a subharmonic function, one variable at a time.
Gauss norms and torus means are multiplicative, in logarithms
additive; the scalar and the power of \(Y_0\) have global
height zero.
For a product of nonzero polynomials the logarithmic torus mean satisfies
\[
  \int_{\mathbb T^n}\log|FG|=\int_{\mathbb T^n}\log|F|+\int_{\mathbb T^n}\log|G| .
\]
 This proves the comparison.

\proofstep{The coefficient estimate when the height is positive.}
There is nothing to prove when \(h_0(F_Z)\le0\). Otherwise
\(H=h_{\mathbb T}(q)>0\). Consider coefficient vectors of
polynomials homogeneous of degree \(h\) in \(Y\) and degree
at most \(h'\) in \(U\). Their dimension and the dimension of
the multiples of \(q\) are, respectively,
\[
 A=\binom{h+s+1}{s+1}\binom{h'+M}{M},\qquad
 B=\binom{h-d_0+s+1}{s+1}\binom{h'-u+M}{M}.
\]
Indeed both degrees add exactly in a product in a polynomial
domain, and the quotient of a homogeneous polynomial by a
homogeneous factor is homogeneous. Thus the shifted monomials
times \(q\) give a basis for the second space. We need
\[
 \frac{B}{A-B}\ge \frac h{(s+2)d_0},\qquad
                    \log A+P_O\le h .
\]

For the first assertion use the product expressions for \(B/A\).
The sum of the deficits of their factors is at most
\[
 (s+1)d_0/h+Mu/h'\le(s+1+1/8)d_0/h.
\]
Since \(d_0/h\le 1/(8M^2(s+1))\), division by one minus
this last deficit gives 
\[
  A/B-1\le(s+2)d_0/h.
\]
For the second assertion bound \(A\le(h+1)^{s+2}(h'+1)^M\)
by counting exponents separately. Now \(h'\le h^2\),
\(h\ge32M^2\) for \(s\ge3\), and \(\log(h+1)\le\sqrt h\).
Hence \(\log A\le3M\sqrt h\le3h/\sqrt{32}\), leaving
more than \(P_O\) by the definition of \(h\).

\proofstep{The metric on the quotient of the coefficient lattice.}
Give the coefficient spaces the Euclidean coefficient norms.
The normalized Arakelov degree of the saturated sublattice of
multiples in \(\mathcal O_K^A\) is at most \(-BH\).
Here is a direct determinant calculation. At a finite place, multiplicativity of the Gauss norm identifies the cofactor lattice by
\[
  \|qg\|_{\rm Gauss}\le1\quad\Longleftrightarrow\quad \|g\|_{\rm Gauss}\le\|q\|_{\rm Gauss}^{-1}.
\]
At an embedding, order the cofactor monomials lexicographically.
Gram--Schmidt on their multiples by \(q\) gives at each step
the norm of \(q\) times a polynomial with leading coefficient
one, measured in the coefficient \(L^2\) norm of the product.
For the cofactor multiple \(H\), Parseval followed by Jensen gives
\[
  \Big(\sum_\alpha |[X^\alpha]H|^2\Big)^{1/2}=\|H\|_{L^2(\mathbb T^n)}\ge\exp\!\left(\int_{\mathbb T^n}\log|H|\right).
\]
 The mean
logarithm of the monic factor is nonnegative by iterated
Jensen. Thus each orthogonal length is at least the
exponential of the mean log absolute value of \(q\). Multiplying
these bounds and inserting the finite Gauss balls gives
the claimed degree with the indicated sign.

The orthogonal quotient lattice therefore has normalized
degree at least \(BH\) in rank \(A-B\). Minkowski's convex
body theorem applied to the lattice of real dimension
\([K:\mathbb Q](A-B)\) supplies an integral coefficient
vector \(Q\) with nonzero image such that its quotient norms
satisfy
\[
 \operatorname{mean}_\sigma\log\|[Q]\|_\sigma
    \le-\frac{BH}{A-B}+\tfrac12\log A+P_O .
\]
In detail, the real covolume in the conjugate-compatible
sum metric is \(|\operatorname{disc} K|^{(A-B)/2}\)
times the exponential of minus the unnormalized degree.
The Euclidean unit ball in dimension \(a\) contains the
cube of side \(2/\sqrt a\). Applying Minkowski and then
arithmetic--geometric mean to the embedding norms cancels
\(\sqrt{[K:\mathbb Q]}\), leaving \(\sqrt{A-B}\)
and half the log root discriminant. The bound above enlarges
both allowances. The quotient is the image of the integral
coefficient lattice, so its nonzero vector lifts to such a
\(Q\).

\proofstep{Applying the integral section bound at the prescribed weight.}
After substituting 
\[
  Y_i=\sum_l U_{il}w_l,
\]
 some coefficient
\(g\) of the resulting polynomial in \(U\) is nonzero on
\(Z\) by the kernel assertion. This is a single integral
combination of \(h\)-fold products. To bound it at an embedding
one may subtract the orthogonal projection of \(Q\) onto the
space of multiples of \(q\) over \(\mathbb C\), since the
substitution of the difference has the same value on \(Z\).
Its coefficient one-norm is at most \(\sqrt A\|[Q]\|_\sigma\).

If \((c_i)_{1\le i\le A}\) are the coefficients of the orthogonal representative, Cauchy--Schwarz gives
\[
  \sum_{i=1}^{A}|c_i|\le\sqrt A\left(\sum_{i=1}^{A}|c_i|^2\right)^{1/2}=\sqrt A\,\|[Q]\|_\sigma .
\]
Each substituted \(Y\)-monomial contributes at most
\(N^h e^{h m_\sigma}\) in the pointwise metric, by the
triangle inequality for the expansion. The same upper bound
holds for every extracted coefficient times the one-norm.
Comparing with \eqref{eq:chow-lower-section} gives
\[
 \frac{hH}{(s+2)d_0}\le
 h(P_0 S_L+G_c S_L+\log N)+\log A+P_O .
\]
The two monomial counts and \(m_Zd_0\le\delta_Z\) prove the
Chow bound for \(F_Z\) and summation proves the required
bound for the family. Although \(h,h'\) grow with the number
of rows and columns, both cancel in this estimate. The
remaining arithmetic input in this step is the componentwise
lower bound at precisely \(h w_L\).
\label{cor:slice-chow-1}
\end{proof}
